% Chapter 7

\chapter{Conclusion} % Main chapter title

\label{Chapter7} % For referencing the chapter elsewhere, use \ref{Chapter7} 

\section{Summary of Empirical Findings}

This thesis has applied a comprehensive macroeconomic framework developed by \textcite{jondeau_environmental_2026} to quantify and compare the global environmental pressures linked to the financial systems of the United States and the Euro area from 2000 to 2022. By integrating national financial accounts with the environmentally extended multi-regional input-output (EE-MRIO) database EXIOBASE v3, the analysis yields three central insights regarding the structural differences, geographical spillovers, and multi-dimensional footprints of major macro-financial systems.

First, the empirical results confirm that structural domestic differences determine the channels through which environmental pressures are financed. In the United States, a fundamentally market-based system, financed emissions are heavily concentrated within investment and pension funds, reflecting their substantial capital allocation toward domestic non-financial corporations (NFCs). Conversely, the Euro area’s bank-based framework routes its environmental impacts primarily through traditional Monetary Financial Institutions (MFIs) and specialized Other Financial Institutions (OFIs) that increasingly finance emitters across borders. 

Second, a stark divergence has emerged in the geographical footprint of the two financial regimes. While the US financial system has consistently maintained an inward focus, allocating over 80\% of its exposure to domestic sectors, the Euro area financial system has increasingly become an engine for cross-border financing. Driven by post-2012 structural dynamics, including the stagnation of internal sovereign debt financing linked to the bank-sovereign "doom loop" and an extended period of negative interest rate policies, Euro area institutions aggressively expanded their foreign claims. By 2022, the Rest of the World (RoW) sector accounted for approximately 40\% of the Euro area's total exposure. 

Third, this outward allocation means that the Euro area has effectively "outsourced" its environmental footprint. When emerging markets and China are incorporated into the investment pool (the RoW3 scenario), the Euro area's financed emissions surpass those of the US in recent years. Remarkably, while absolute financed emissions in the US represent a smaller fraction of its massive domestic physical profile, the Euro area's financed emissions range is roughly equivalent to its entire territorial. This global footprint is not limited to carbon; the Euro area financial system exhibits an equally disproportionate external dependence on global ecosystems via cross-border claims linked to intensive land use, blue-water depletion, and upstream mineral and fossil fuel extraction.

\section{Policy Implications for Central Banks and Regulators}

The empirical findings of this thesis carry several implications for central banks, macro-prudential supervisors, and international financial standard-setters navigating the climate transition.

First, the results suggest that financial regulation can shape financed emissions in unintended ways by influencing the allocation of capital across sectors and regions. For example, Basel III liquidity rules, such as the Liquidity Coverage Ratio, may help anchor US financial capital in relatively low-carbon Treasury securities, while Euro area sovereign-risk dynamics and yield-search incentives may encourage greater exposure to more carbon-intensive foreign assets. This implies that financial stability policies should not be assessed solely through the lens of domestic risk reduction, but also in relation to their broader environmental footprint.

Second, the findings have implications for climate stress testing. Existing exercises, such as the ECB's 2022 climate risk stress test and the Federal Reserve's pilot Climate Scenario Analysis exercise, provide an important starting point for assessing climate-related financial risks at the bank level. However, because these frameworks are mainly designed around banks' exposures to clients, counterparties, and selected credit portfolios, they may be less suited to capturing macro-financial risks arising from the aggregate allocation of capital across sectors and countries. This thesis therefore highlights the value of complementing existing supervisory exercises with system-level, multi-regional approaches that account for cross-border financial intermediation and the carbon intensity of foreign assets.

Third, by extending the macro-financial framework beyond greenhouse gas emissions to land use, water consumption, and resource extraction, this thesis speaks to the broader principle of double materiality. Environmental risks are not limited to carbon emissions: financial systems may also be exposed to physical vulnerabilities arising from water scarcity, land degradation, biodiversity loss, or resource dependence. This is particularly relevant for highly internationalized financial systems such as the Euro area, where cross-border claims may finance production networks with substantial non-carbon environmental pressures. These findings suggest that climate-related financial supervision should gradually move beyond carbon-focused metrics and incorporate a wider set of environmental indicators. More harmonized disclosure and satellite-accounting frameworks could help supervisors and financial institutions identify portfolios that appear relatively climate-aligned in terms of emissions, but remain exposed to significant environmental pressures through land, water, or resource channels.

\section{Future Research and Concluding Remarks}

One of the main assumptions of the framework is that the geographic distribution of final borrowers is not directly observed. As a result, the model assumes that financial systems fund developed economies, emerging economies, or both groups proportionally, based on each group’s claims relative to total global claims. If more granular information on the national profile of credit recipients becomes available, future research could incorporate this information to improve the accuracy of the attribution exercise.

Second, the current public coverage of EXIOBASE extends only to 2022, and the most recent years should be interpreted with caution. This highlights the need for the continued development of environmentally extended multi-regional input-output databases that combine broad country and sectoral coverage with frequent updates and detailed environmental indicators. Future work could therefore contribute to the construction of reliable and regularly updated EE-MRIO sources, particularly for applications in climate finance and macro-financial risk assessment.

Third, future research could reconcile the macroeconomic attribution framework used in this thesis with bottom-up financed-emissions measures reported by major financial institutions. Such a comparison would help assess the completeness and consistency of existing emissions and environmental-pressure disclosures, while also identifying potential gaps across institutions, asset classes, and jurisdictions.

In summary, as ecological and climate crises accelerate, the definition of a “safe” or “sustainable” financial system must undergo a fundamental shift. This thesis shows that an economy cannot achieve a genuine green transition through the decarbonization of domestic production alone if its financial intermediaries continue to enable environmental degradation across borders. Measuring sustainability purely through territorial production or final consumption metrics overlooks the structural role of globalized financial capital in shaping environmental outcomes. By adopting a comprehensive financial-environmental attribution framework, central banks, regulators, and market participants can more accurately evaluate, disclose, and ultimately mitigate the global ecosystem pressures enabled by modern financial systems.
